% !TeX program = xelatex
% =============================================================================
% 中国发明专利申请文件（草稿）——LaTeX 模板
% =============================================================================
%
% 编译链（四步，缺一不可）：
%   xelatex -interaction=nonstopmode main.tex
%   bibtex  main
%   xelatex -interaction=nonstopmode main.tex
%   xelatex -interaction=nonstopmode main.tex
%
% 为什么必须 XeLaTeX：中文要调用系统字体，pdfLaTeX 编 ctex 会直接报错。
% 为什么必须跑两遍收尾：参考文献、交叉引用、页码都要第二遍才收敛。
%
% 文件顺序按《专利法实施细则》对申请文件的要求排：说明书摘要 →
% 权利要求书 → 说明书。请求书是国知局表格，不在本模板内。
%
% 排版参数（页边距/字号/行距）集中在下面的 geometry 与各 \zihao 处。
% 纸张、页边、字体、字高、行距、页码的**强制性依据**在《专利审查指南》第五部分
% 第一章 4.1–5.6；本模板的 geometry 与它们一一对应。
% 强制性要求（纸张、字迹、各部分标题与顺序、摘要字数、附图规范、电子申请格式）
% 见 references/format.md。正式提交前请用 CNIPA 官方渠道核对，并建议交由执业
% 代理师复核。
% 注：模板里的途径引用（\cite）只是演示，正式提交时引证的对比文件应当写全
% 公开号/文献出处；请勿把 refs.bib 里的示例条目原样保留。
%
% ★ 电子申请的一个坑：官方《关于规范提交专利电子申请的指引》明确要求
%   「说明书不应添加任何形式的段落编号」——新申请 XML 的说明书段号由系统自动
%   生成。本模板的 [0001] 段号是**草稿阶段**方便按段答复审查意见用的，正式提交
%   前必须关掉：把下面的 \paranumbertrue 改成 \paranumberfalse。
%
% 模板里的技术内容全部是**演示用示例**（一种"示例装置"），不是真实申请。
% =============================================================================

\documentclass[12pt,a4paper,fontset=windows]{ctexart}

% ---- 版式 -------------------------------------------------------------------
\usepackage[top=25mm,bottom=15mm,left=25mm,right=15mm,headsep=5mm,footskip=8mm]{geometry}
\usepackage{fancyhdr}
\usepackage{lastpage}

% ---- 数学、表格、图形 -------------------------------------------------------
\usepackage{amsmath,amssymb}
\usepackage{booktabs}
\usepackage{array}
\usepackage{graphicx}
\usepackage{enumitem}
\usepackage[hidelinks]{hyperref}

% ---- 目录：CNIPA 申请文件不要求目录，草稿阶段可打开方便自查 -----------------
% 正式提交前务必保持关闭（\tableofcontentsfalse）。
\newif\ifshowtoc
\showtocfalse

% ---- 段落编号 [0001] --------------------------------------------------------
% 说明书段落编号不是提交时必须的，但**答复审查意见时按段引用**极为方便，
% 也是公开文本（CN…A）的编号方式。
% 开关：草稿默认开；正式提交（尤其电子申请）务必改成 \paranumberfalse——
% 官方指引要求说明书不应添加任何形式的段落编号，段号由系统自动生成。
\newif\ifparanumber
\paranumbertrue
\newcounter{specpara}
\newcommand{\padfour}[1]{%
  \ifnum#1<10 [000#1]\else
  \ifnum#1<100 [00#1]\else
  \ifnum#1<1000 [0#1]\else
  [#1]\fi\fi\fi}
% 注意：这里必须用 \thespecpara，不能写 \value{specpara}——
% \ifnum 后面跟 \value{...} 会报 "Missing number, treated as zero"（已实测）。
\newcommand{\spara}[1]{\par\refstepcounter{specpara}%
  \ifparanumber\noindent\padfour{\thespecpara}\ \fi
  #1\par}
% 不带编号的段落（用于摘要、发明内容里的小标题下正文等）
\newcommand{\npara}[1]{\par\noindent #1\par}

% ---- 结构宏 ----------------------------------------------------------------
% 一级部件标题（说明书摘要 / 权利要求书 / 说明书）：三号黑体居中
\newcommand{\parttitle}[1]{%
  \par\vspace{0.6em}%
  \begin{center}{\zihao{3}\heiti #1}\end{center}%
  \vspace{0.3em}\par}
% 说明书内部的法定小节标题（技术领域 / 背景技术 / ……）：小四黑体顶格
\newcommand{\sect}[1]{\par\vspace{0.5em}\noindent{\zihao{-4}\heiti #1}\par\vspace{0.2em}}
% 权利要求：编号 + 正文，编号顶格
\newcommand{\claim}[2]{\par\noindent #1. #2\par}
% 附图标记：正文里首次出现用括号标注，全文必须与附图一致
\newcommand{\mref}[1]{（#1）}
% 引用对比文件
\newcommand{\cmpfile}[1]{\cite{#1}}

% ---- 页眉页脚 --------------------------------------------------------------
% 申请文件正文一般不设页眉；页码保留在页脚居中，便于校对与答复时指页。
\pagestyle{fancy}
\fancyhf{}
\renewcommand{\headrulewidth}{0pt}
\cfoot{\small 第 \thepage\ 页\quad 共 \pageref{LastPage} 页}

% ---- 元信息（草稿用，正式提交的请求书另外填写） -----------------------------
\newcommand{\drafttitle}{一种示例装置}
\newcommand{\draftapplicant}{申请人名称（示例）}
\newcommand{\draftdocket}{代理机构案卷号（示例）：CN-0000-000}

\title{\drafttitle}
\author{\draftapplicant}
\date{\today}

% =============================================================================
\begin{document}

% 草稿题头：标明这是工作稿，正式提交时删掉这一段
\begin{center}
  {\zihao{-4}\heiti 发明专利申请文件（草稿）}\\[0.3em]
  {\small \draftdocket}
\end{center}
\vspace{0.5em}

\ifshowtoc
  \tableofcontents
  \newpage
\fi

% -----------------------------------------------------------------------------
\parttitle{说明书摘要}
% -----------------------------------------------------------------------------
% 摘要：写明所解决的技术问题、技术方案要点与主要用途，不得使用商业性宣传
% 用语，不得写"如权利要求……所述"这类引用语。字数上限见 references/format.md。
\npara{本发明公开了一种示例装置，属于机械传动技术领域。所述装置包括壳体、
设于所述壳体内的驱动单元以及由所述驱动单元驱动的输出部件；所述驱动单元与
所述输出部件之间设有离合组件，所述离合组件被配置为在所述输出部件的负载超过
预设阈值时断开动力传递。本发明通过上述离合组件，解决了现有装置在堵转工况下
易损坏驱动单元的问题，提高了装置的可靠性与使用寿命。本发明可用于需要频繁
启停的自动化设备。}

\vspace{0.5em}
\noindent{\zihao{-4}\heiti 摘要附图：}图 1。

% -----------------------------------------------------------------------------
\parttitle{权利要求书}
% -----------------------------------------------------------------------------
% 权利要求撰写规则、独立/从属权利要求的格式要求见 references/claims.md。
% 要点：编号从 1 起连续；独立权利要求用"其特征在于"划分前序部分与特征部分；
% 从属权利要求必须引用在前的权利要求，并写明附加技术特征。

\claim{1}{一种示例装置，其特征在于，包括：
  壳体\mref{1}；
  驱动单元\mref{2}，设于所述壳体\mref{1}内；以及
  输出部件\mref{3}，由所述驱动单元\mref{2}驱动；
  其中，所述驱动单元\mref{2}与所述输出部件\mref{3}之间设有离合组件\mref{4}，
  所述离合组件\mref{4}被配置为在所述输出部件\mref{3}的负载超过预设阈值时
  断开动力传递。}

\claim{2}{根据权利要求 1 所述的示例装置，其特征在于，所述离合组件\mref{4}包括
  摩擦片\mref{41}和弹性件\mref{42}，所述弹性件\mref{42}将所述摩擦片\mref{41}
  压紧于所述驱动单元\mref{2}的输出端。}

\claim{3}{根据权利要求 2 所述的示例装置，其特征在于，所述弹性件\mref{42}为
  螺旋弹簧或碟形弹簧。}

\claim{4}{根据权利要求 1 至 3 中任一项所述的示例装置，其特征在于，所述壳体%
  \mref{1}上设有散热孔\mref{11}。}

\claim{5}{一种自动化设备，其特征在于，包括权利要求 1 至 4 中任一项所述的
  示例装置。}

% -----------------------------------------------------------------------------
\parttitle{说明书}
% -----------------------------------------------------------------------------
% 指南第一部分第一章 4.2：说明书第一页第一行应当写明发明名称，与请求书一致，
% 左右居中，前面不得冠以「发明名称」或「名称」等字样，与正文之间空一行。
% 名称一般不超过 25 字，必要时不超过 60 字（指南第二部分第二章 2.2.1）。
\begin{center}{\zihao{-3}\heiti \drafttitle}\end{center}
\vspace{0.5em}

\sect{技术领域}
\spara{本发明涉及机械传动技术领域，具体涉及一种示例装置及包括该装置的自动化设备。}

\sect{背景技术}
% 背景技术：写明对发明的理解、检索、审查有用的背景技术，并引证反映这些
% 背景技术的文件；只描述现有技术的问题，不要在这里写本发明的方案。
\spara{现有的传动装置通常将驱动单元与输出部件直接连接。例如，\cmpfile{examplepatent}%
公开了一种将电机输出轴与负载直接连接的传动结构；\cmpfile{examplepaper}指出，
该类结构在负载突变时容易产生冲击。}
\spara{上述现有技术存在如下问题：当输出部件因异物卡阻而堵转时，驱动单元仍需
持续输出扭矩，容易导致驱动单元过热甚至烧毁；同时，直接连接的结构无法在堵转
工况下自动切断动力，设备的维护成本较高。}

\sect{发明内容}
% 发明内容：写明要解决的技术问题、技术方案、以及相对于现有技术的有益效果。
% 技术问题的表述应与权利要求的技术特征对应，不要引入权利要求之外的新特征。
\spara{针对现有技术存在的上述问题，本发明要解决的技术问题是：如何在输出部件
负载异常时自动切断动力传递，以避免驱动单元损坏。}
\spara{为解决上述技术问题，本发明提供一种示例装置，包括壳体、设于所述壳体内的
驱动单元、以及由所述驱动单元驱动的输出部件；所述驱动单元与所述输出部件之间
设有离合组件，所述离合组件被配置为在所述输出部件的负载超过预设阈值时断开
动力传递。}
\spara{与现有技术相比，本发明具有以下有益效果：其一，堵转工况下动力被自动
切断，驱动单元不再持续过载；其二，负载恢复正常后装置可继续工作，无需更换
驱动单元；其三，结构简单、成本低，便于在既有设备上改造实施。}

\sect{附图说明}
\spara{图 1 为本发明实施例提供的示例装置的整体结构示意图。}
\spara{图 2 为本发明实施例提供的离合组件的局部放大示意图。}
\spara{附图标记说明：1—壳体；11—散热孔；2—驱动单元；3—输出部件；
4—离合组件；41—摩擦片；42—弹性件。}

\sect{具体实施方式}
% 具体实施方式：至少要给出一个能实现该发明的具体技术方案；有附图的结合附图
% 说明；有数值/范围的应给出；不要只重复权利要求书的文字。
\spara{下面结合附图和具体实施例对本发明作进一步说明。所述实施例仅用于说明
本发明的技术方案，不用于限制本发明的保护范围。}
\spara{参见图 1，本发明实施例提供的示例装置包括壳体 1、驱动单元 2、输出部件 3
以及离合组件 4。壳体 1 内部形成容纳空间，驱动单元 2 通过螺栓固定于壳体 1 的
底壁。在本实施例中，驱动单元 2 采用额定功率 200 W 的直流无刷电机，额定转速
为 3000 r/min；在其它实施例中，也可以采用伺服电机或气动马达，本发明对此
不作限定。}
\spara{参见图 2，离合组件 4 设于驱动单元 2 的输出端与输出部件 3 之间，包括
摩擦片 41 和弹性件 42。弹性件 42 的一端抵靠壳体 1，另一端将摩擦片 41 压紧于
驱动单元 2 的输出端。在本实施例中，弹性件 42 为螺旋弹簧，其预紧力为 15 N；
当输出部件 3 承受的负载扭矩超过由该预紧力确定的阈值时，摩擦片 41 与驱动单元 2
的输出端之间产生相对滑动，从而断开动力传递。}
\spara{在本实施例中，壳体 1 的侧壁上开设有散热孔 11，用于在装置正常运行时
将驱动单元 2 产生的热量排出。散热孔 11 的数量为 6 个，沿壳体 1 的周向均匀
分布；在其它实施例中，散热孔 11 的数量和布置方式可根据散热需求调整。}
\spara{本实施例的示例装置可应用于自动化设备。将上述示例装置安装于自动化设备
的进料机构中，当进料机构被异物卡阻时，离合组件 4 自动断开动力传递，避免
驱动单元 2 损坏；异物清除后，摩擦片 41 在弹性件 42 的作用下重新压紧，装置
恢复正常工作。}
\spara{以上所述仅为本发明的较佳实施例，并不用于限制本发明。凡在本发明的
精神和原则之内所作的任何修改、等同替换和改进等，均应包含在本发明的保护
范围之内。}

% -----------------------------------------------------------------------------
% 说明书附图（正式提交时通常另页绘制；此处给出占位说明）
% -----------------------------------------------------------------------------
\newpage
\parttitle{说明书附图}
\noindent 图 1　示例装置整体结构示意图\par
\vspace{0.5em}
\begin{center}
  \fbox{\parbox[c][35mm][c]{0.8\textwidth}{\centering\small
    此处插入图 1：用绘图软件绘制后以 PDF/EPS 矢量图插入。\\
    附图要求（线条、编号、文字限制）见 references/format.md。}}
\end{center}
\vspace{1em}
\noindent 图 2　离合组件局部放大示意图\par
\vspace{0.5em}
\begin{center}
  \fbox{\parbox[c][35mm][c]{0.8\textwidth}{\centering\small
    此处插入图 2。附图标记须与说明书正文中的标记一致。}}
\end{center}

% -----------------------------------------------------------------------------
% 引证的对比文件 / 参考文献
% -----------------------------------------------------------------------------
% 背景技术中引证的对比文件必须在此可查；正文用 \cmpfile{键名} 引用。
% 用 unsrt 而不是 plain：对比文件按正文引用顺序编号（[1]、[2]…）更符合阅读习惯。
\newpage
\renewcommand{\refname}{引证文件}
\bibliographystyle{unsrt}
\bibliography{refs}

\end{document}
